\IEEEPARstart{W}{e} consider the problem of defect detection in medical devices, which is a critical task to ensure the safety and reliability of these devices. Such defects can arise due to various reasons, including hardware malfunctions, software bugs, or environmental factors. Detecting these anomalies promptly is essential to prevent potential harm to patients and to maintain the integrity of the medical device. \emph{The complexity of modern medical devices, which often incorporate advanced semiconductor components, has made traditional defect detection methods less effective.} As a result, there is a growing interest in leveraging machine learning techniques to enhance the detection process and organize data in a more structured manner using knowledge graphs (cf.\cite{lee2017convolutional,khader2018stencil,lu2019realtime,cao2024xscan}.

In particular, \cite{shenoy2026uncertainty} introduces a novel approach to defect–mechanism–parameter violation hypotheses for surface-mount assembly root-cause analysis from uncertain observations. Using OKN as the source of information, the authors propose a method to generate hypotheses that can explain the observed defects. In doing so, they leverage stage-wise three-head \gls*{mlp} architecture that models data uncertainty by estimating defect and mechanism likelihoods and parameter-level risk scores, enabling probabilistic and traceable root-cause ranking. Thus, they are able to detect defects (head 1), identify the underlying mechanisms (head 2), and assess the associated risks (head 3). However, this approach is not without its limitations. In particular, the authors note that their method can be sensitive to epistemic uncertainty, which can lead to detection failures and misclassification of defects. This is a significant concern in high-stakes applications such as medical device manufacturing, where misclassification can incur substantial costs and pose a significant risk of serious patient harm.
\begin{enumerate}
    \item \marco{Mention the work that has inspired this one and the gap that we are trying to fill}
    \item \marco{mention and explain the role of the heads, crucailly if head 1 fails, the rest of the infomration is worthless}
    \item \marco{Explain epistemic uncertainty and its relevance to the problem}
\end{enumerate}

\subsection{Contributions}
\label{subsec:contributions}
We summarize the contributions of this work as follows:
\begin{enumerate}
    \item We investigate detection failures arising from epistemic uncertainty in recent root-cause analysis approaches for semiconductor manufacturing, a high-stakes application in which misclassification can incur substantial costs and, in the context of medical device production, pose a significant risk of serious patient harm.
    \item \marco{We analyse a simple, yet effective, principled approach to mitigate the impact of epistemic uncertainty on detection failures by leveraging a cost-sensitive classification framework that incorporates uncertainty handling into the decision-making training process. cite prior papers makig sure to say that one is the inspiration, and the other required a more principled analysis, which allows to find }
    \item We showcase the effectiveness of our approach through extensive experiments on syntetic benchamrk
    \item We consider a real-world case study of a medical device, introducing a benchmar easily downloadable through Proto-OKN, demonstrating the practical applicability and benefits of our proposed method in a high-stakes setting.
\end{enumerate}